\chapter{Conclusão}
\label{cha:conclusao}
\index{Conclusão}

Este estudo investiga, em caráter exploratório, as técnicas de
reamostragem \textit{bootstrap}, \textit{bagging} e \textit{subagging}
aplicadas à Regressão Linear Múltipla (RLM) na calibração multivariada
de amostras de aço-ligas. O trabalho implementa as abordagens
utilizando a RLM como modelo base e atende aos objetivos propostos,
dentro das limitações discutidas ao final desta seção.

%Os experimentos seguem um protocolo controlado, com a otimização dos hiperparâmetros por validação cruzada 5-\textit{fold} aplicada exclusivamente ao conjunto de calibração. Esse protocolo busca assegurar comparações justas e reprodutíveis entre a \textit{RLM clássica}, o \textit{RLM-bagging} e o \textit{RLM-subagging}. Além do RMSEP, o estudo incorpora as métricas \gls{MAE} e R$^2$, ampliando acaracterização do desempenho preditivo de cada técnica. Ressalta-se que as conclusões a seguir se referem a esse protocolo com otimização por validação cruzada: nos experimentos anteriores, com parâmetros fixados sem otimização formal, a técnica vencedora em alguns analitos (Ni e Fe) foi outra. Isso reforça que o desempenho relativo depende da forma como $m$ e $k$ são selecionados.

Os experimentos seguem um protocolo controlado, com a otimização dos
hiperparâmetros por validação cruzada 5-\textit{fold} aplicada
exclusivamente ao conjunto de calibração. Esse protocolo busca
assegurar comparações justas e reprodutíveis entre a \textit{RLM
clássica}, o \textit{RLM-bagging} e o \textit{RLM-subagging}. Além do
RMSEP, o estudo incorpora as métricas \gls{MAE} e R$^2$, ampliando a
caracterização do desempenho preditivo de cada técnica.

Ressalta-se que as conclusões a seguir se referem a esse protocolo com 
otimização por validação cruzada: nos experimentos anteriores, com 
parâmetros fixados sem otimização formal, a técnica vencedora em alguns 
analitos (Ni e Fe) foi outra. Isso reforça que o desempenho relativo 
depende da forma como $m$ e $k$ são selecionados.

%----------------------------------------------------------------------------

A análise comparativa revela variações de desempenho dependentes do
analito e indica, para este conjunto de dados específico, a
inexistência de uma técnica universalmente superior. Para o Manganês
(Mn), o RLM-\textit{subagging} ($k^* = 80\%$, $m^* = 10$) apresenta o
menor RMSEP (0,0966), uma redução de 8,1\% em relação à RLM clássica.
Para o Ferro (Fe), o RLM-\textit{bagging} ($m^* = 30$) obtém o melhor
desempenho (RMSEP = 2,6406), com melhoria de 4,0\% sobre a RLM
clássica.

Para o Níquel (Ni), o RLM-\textit{subagging} ($k^* = 60\%$,
$m^* = 10$) exibe resultado marginalmente superior (RMSEP = 1,6003),
com diferença de 0,25\% em relação à RLM clássica. Dada a magnitude
dessa diferença e o tamanho amostral disponível, o resultado deve ser
interpretado com cautela. Em contrapartida, para o Molibdênio (Mo) e o
Cromo (Cr), a RLM clássica apresenta desempenho superior (RMSEP de
0,2633 e 1,0798, respectivamente), indicando que a complexidade
adicional das estratégias de \textit{ensemble} não resulta em ganho de
precisão para esses analitos neste contexto.

%---------------------------------------------------------------------

A avaliação por meio do RMSEP, complementada pelo \gls{MAE} e pelo
R$^2$, permite uma caracterização abrangente do desempenho preditivo.
Todos os modelos apresentam coeficiente de determinação superior a
0,98 para todos os elementos, o que evidencia elevada capacidade
preditiva para o conjunto de dados analisado.

Para Mn e Ni, o RLM-\textit{subagging} apresenta dominância de Pareto
sobre o RLM-\textit{bagging}, por ser simultaneamente superior em
acurácia (RMSEP e \gls{MAE}) e mais eficiente computacionalmente, o
que reforça sua indicação nesses dois casos. Para Mo e Cr, os
resultados sustentam a adoção da RLM clássica, em razão de sua
simplicidade e de sua eficiência computacional.

%--------------------------------------------------------------------

A análise da eficiência computacional mostra diferenças expressivas
entre os métodos. A RLM clássica é a abordagem mais rápida em todos
os analitos. O RLM-\textit{subagging} ocupa posição intermediária:
tomando a razão entre os tempos do RLM-\textit{bagging} e do
RLM-\textit{subagging} em cada analito, o subagging é, em média, 3,0
vezes mais rápido, com variação entre 1,4 vez (Mo) e 6,3 vezes (Ni).
Em termos agregados, a soma dos tempos dos cinco analitos é cerca de
53\% menor no subagging.

Em relação à RLM clássica, o RLM-\textit{subagging} apresenta custo de
até 19,5 vezes superior, enquanto o RLM-\textit{bagging} é, em média,
22,1 vezes mais lento. O Cromo (Cr) é o único caso em que o
RLM-\textit{subagging} não apresenta vantagem temporal sobre o
RLM-\textit{bagging}, comportamento explicado pelo maior número de
replicações selecionado por validação cruzada ($m^* = 30$ contra
$m^* = 20$).

%-------------------------------------------------------------------------

A investigação indica que o RLM-\textit{subagging} é a alternativa
mais equilibrada entre precisão e eficiência para os dados analisados,
sendo recomendado para Mn e Ni. Para Mo e Cr, os resultados sustentam
a adoção da RLM clássica. O RLM-\textit{bagging} mostra-se preferível
para Fe, embora o RLM-\textit{subagging} atinja desempenho preditivo
equivalente (RMSEP de 2,6509 e R$^2$ idêntico, 0,9900) em menos da
metade do tempo. A relação de compromisso entre acurácia e tempo atua
como elemento decisivo na seleção da técnica.

% Do ponto de vista da engenharia de software, o trabalho atinge seus objetivos específicos. A arquitetura modular em três camadas independentes (\texttt{models.py}, \texttt{components.py} e \texttt{app.py}) satisfaz o requisito de modularidade (RNF01). A suíte de 92 testes automatizados atesta a equivalência entre o sistema Python e o sistema Java legado (RNF03) em dois níveis distintos: para a RLM, que é determinística, a equivalência numérica é estrita (tolerância de $5\times10^{-4}$ em RMSEP, \gls{MAE} e R$^2$); para o \textit{bagging} e o \textit{subagging}, que são estocásticos e cujas sequências pseudoaleatórias diferem entre Python e Java, a suíte verifica consistência estatística dentro de tolerâncias mais amplas (0,15 em RMSEP e \gls{MAE}; $5\times10^{-3}$ em R$^2$). Assim, a migração incremental orientada por testes de caixa-preta constitui contribuição metodológica válida independentemente do domínio de aplicação, ainda que não demonstre identidade das reamostragens entre as duas implementações.

Do ponto de vista da engenharia de software, o trabalho atinge seus
objetivos específicos. A arquitetura modular em três camadas
independentes (\texttt{models.py}, \texttt{components.py} e
\texttt{app.py}) satisfaz o requisito de modularidade (RNF01). A suíte
de 92 testes automatizados atesta a equivalência entre o sistema
Python e o sistema Java legado (RNF03) em dois níveis distintos.

Para a RLM, que é determinística, a equivalência numérica é estrita,
com tolerância de $5\times10^{-4}$ em RMSEP, \gls{MAE} e R$^2$. Para o
\textit{bagging} e o \textit{subagging}, que são estocásticos e cujas
sequências pseudoaleatórias diferem entre Python e Java, a suíte
verifica consistência estatística dentro de tolerâncias mais amplas:
0,15 em RMSEP e \gls{MAE}; $5\times10^{-3}$ em R$^2$.

Assim, a migração incremental orientada por testes de caixa-preta
constitui contribuição metodológica válida independentemente do
domínio de aplicação. Entretanto, essa abordagem não demonstra
identidade das reamostragens entre as duas implementações.

%Este estudo apresenta limitações pertinentes. O conjunto de dados corresponde a uma composição específica de aço-liga, analisada sob condições instrumentais particulares, o que restringe a extrapolação dos achados para outras composições, regimes amostrais ou configurações espectrométricas. A razão de aproximadamente 42 preditores por observação no conjunto de calibração (723 variáveis para 17 amostras) caracteriza um regime de alta dimensionalidade que, embora tratável pela pseudo-inversa de Moore-Penrose, amplifica a sensibilidade dos resultados à composição das amostras disponíveis. Os \textit{folds} da validação cruzada contêm apenas três ou quatro amostras, o que limita a estabilidade da seleção de $m^*$ e $k^*$.

Este estudo apresenta limitações pertinentes. O conjunto de dados
corresponde a uma composição específica de aço-liga, analisada sob
condições instrumentais particulares, o que restringe a extrapolação
dos achados para outras composições, regimes amostrais ou
configurações espectrométricas.

A razão de aproximadamente 42 preditores por observação no conjunto
de calibração (723 variáveis para 17 amostras) caracteriza um regime
de alta dimensionalidade. Embora tratável pela pseudo-inversa de
Moore-Penrose, esse cenário amplifica a sensibilidade dos resultados
à composição das amostras disponíveis.

Além disso, os \textit{folds} da validação cruzada contêm apenas três
ou quatro amostras, o que limita a estabilidade da seleção de $m^*$ e
$k^*$.

% Outras limitações são de natureza metodológica. O valor $k^* = 80\%$ coincide com o limite superior da grade avaliada ($k \in \{0{,}50;\ 0{,}60;\ 0{,}70;\ 0{,}80\}$) em quatro dos cinco analitos, de modo que o ótimo real pode situar-se em frações maiores. As tolerâncias adotadas nos testes de caixa-preta dos métodos estocásticos são amplas em relação às diferenças de desempenho entre as técnicas, de modo que a suíte assegura consistência estatística, mas não discrimina implementações com desempenhos próximos. Por fim, a análise adota exclusivamente a RLM como modelo base.

Outras limitações são de natureza metodológica. O valor $k^* = 80\%$
coincide com o limite superior da grade avaliada
($k \in \{0{,}50;\ 0{,}60;\ 0{,}70;\ 0{,}80\}$) em quatro dos cinco
analitos, de modo que o ótimo real pode situar-se em frações maiores.

As tolerâncias adotadas nos testes de caixa-preta dos métodos
estocásticos são amplas em relação às diferenças de desempenho entre
as técnicas. Desse modo, a suíte assegura consistência estatística,
mas não discrimina implementações com desempenhos próximos.

Por fim, a análise adota exclusivamente a RLM como modelo base.


% Como direcionamento para investigações futuras, este estudo recomenda a expansão do escopo para conjuntos amostrais maiores e diferentes classes de materiais, a ampliação da grade de $k$ além de 80\% e a comparação com modelos preditivos alternativos, como \gls{PLS}, \gls{LASSO} e redes neurais, que permitirão avaliar se os padrões observados se sustentam em regimes mais generalizáveis. Recomenda-se ainda o uso de uma mesma sequência de reamostragens (por exemplo, índices exportados pelo Java e lidos pelo Python) para permitir a comparação amostra a amostra dos métodos estocásticos, e a avaliação de técnicas de \textit{ensemble} mais sofisticadas, como \textit{boosting} e \textit{stacking}.

Como direcionamento para investigações futuras, este estudo recomenda
a expansão do escopo para conjuntos amostrais maiores e diferentes
classes de materiais, a ampliação da grade de $k$ além de 80\% e a
comparação com modelos preditivos alternativos, como \gls{PLS},
\gls{LASSO} e redes neurais. Essas abordagens permitirão avaliar se os
padrões observados se sustentam em regimes mais generalizáveis.

Recomenda-se ainda o uso de uma mesma sequência de reamostragens
(por exemplo, índices exportados pelo Java e lidos pelo Python) para
permitir a comparação amostra a amostra dos métodos estocásticos.
Sugere-se também a avaliação de técnicas de \textit{ensemble} mais
sofisticadas, como \textit{boosting} e \textit{stacking}.



